\begin{statementbox}
\textbf{条件与对应。} 椭圆 $x^2/a^2+y^2/b^2=1$，$a>b>0$，
$c=\sqrt{a^2-b^2}>0$，$e=c/a\in(0,1)$。
左右焦点为 $F_L(-c,0),F_R(c,0)$，对应准线分别为
$\ell_L:x=-a^2/c$、$\ell_R:x=a^2/c$。
对椭圆上的任意点 $P$，
\[
\boxed{\frac{PF_R}{d(P,\ell_R)}=\frac{PF_L}{d(P,\ell_L)}=e}.
\]
距离是到对应准线的垂距。因为 $|x|\le a<a^2/c$，两准距均为正。

\textbf{第二定义的反向。} 对上述任一固定焦点与对应准线，满足该距离比等于 $e$ 的所有点恰构成整条椭圆。
一般地，给定 $F\notin\ell$ 及 $0<e<1$，点集 $PF=e\,d(P,\ell)$ 为椭圆。
焦点在 $y$ 轴时交换坐标，上焦点配上准线，下焦点配下准线。

\medskip
\textbf{为什么需要焦点--准线定义？} 不只是为了给椭圆多一个公式。
把固定点 $F$、固定直线 $\ell$ 及距离条件 $PF=e\,d(P,\ell)$ 放在一起，
只改变常数 $e$：$0<e<1$ 得椭圆，$e=1$ 得抛物线，$e>1$ 得双曲线。
这使三种圆锥曲线拥有\emph{同一个定义框架}；本条 C039 是其中 $0<e<1$ 的情形。

\medskip
\textbf{不背公式的等价写法。} 设 $P(x,y)$ 在这条椭圆上。因为椭圆在两条准线之间，
\[
\begin{aligned}
d(P,\ell_R)&=\frac{a^2}{c}-x=\frac ae-x, &PF_R&=a-ex,\\
d(P,\ell_L)&=\frac{a^2}{c}+x=\frac ae+x, &PF_L&=a+ex.
\end{aligned}
\]
\textbf{理解关键词：} 右边配右边，左边配左边；\emph{焦半径连接点与焦点，准距是垂直到准线的距离}。
准线距中心的距离是 $a/e$，不是 $ae$；$e<1$，所以 $a/e>a$，准线总在椭圆外。
上面两条焦半径的线性式只对椭圆上的点成立，不能随意套给平面上任意点。

\end{statementbox}
